\documentclass[man]{apa7}
\usepackage[american]{babel}
\usepackage{csquotes}
\usepackage[style=apa, backend=biber]{biblatex}
\addbibresource{references.bib}

\usepackage{hyperref}
\hypersetup{colorlinks=true, linkcolor=blue, citecolor=blue}

\usepackage{amsmath, amssymb, graphicx}
\usepackage{listings}
\usepackage{xcolor}

\DeclareRobustCommand{\code}[1]{%
	\begingroup
	\setlength{\fboxsep}{2pt}%
	\colorbox{gray!15}{\ttfamily\small #1}%
	\endgroup
}

\shorttitle{Transmission Electron Microscopy Analysis}
\begin{document}
	\title{Vector Dimensional Analysis of TEM in Racket}
	\author{Andrew D. France}
	\affiliation{Tarrant County College}
	\keywords{Scheme, Bragg's Law, Trigonometry, TEM, Biotech, Mathematics, Racket, Computer Science, Vector Dimensional Analysis, Case Study, Computational Chemistry, Molecular Biology}
	\abstract{
	A novel method of vector dimensional analysis is implemented in a dialect of \textit{Scheme} called \textit{Racket} as a proof of concept. Prioritizing unit-first algorithms over traditional formula-first numerical methods by treating geometric relationships, such as Bragg's Law, as structural invariants is more mathematically reliable than the alternative. The method of vector dimensional analysis provides a sanity check embedded directly within the framework itself, rather than as a numerical calculation done post-hoc.
	}
	
	\maketitle
	
	\section{Introduction}
	The Biotechnology industry needs reliable methods for computationally calculating the distances between atoms in biological macromolecules.
	Structures like the 1GZX oxy T-state haemoglobin \parencite{1gzx} exist in the Protein Data Bank with documented interatomic coordinates. Traditional computational methods for geometric analysis of these biological structures use formula-based algorithms that can, in theory, introduce dimensional errors. Amit Singer and Fred J. Sigworth describe the goal of structural biology in terms of these structures: ``...to determine their 3D structure—that is, the position of each atom in the machine—in order to try to understand the way in which the machine works.'' \parencite{Singer2020}.
	
	Transmission electron microscopy (TEM) relies on electron diffraction to determine atomic coordinates, with measurements that must satisfy the geometric constraints defined by Bragg's Law. However, traditional formula-first algorithms lack structural invariants by treating these relationships as naked numbers; consequently hiding dimensional assumptions that can lead to drastic computational errors.
	
	Drawing a structural analogy from materials science, Stephen W. Tsai and Jose Daniel D. Melo demonstrated that in the design of laminates, invariants, such as the trace of a stiffness matrix, exist as material properties independent of ply orientation: ``The concept of invariance was proven useful in the design of laminates because the invariants are not affected by ply orientation.'' \parencite{2014Tsai}. \linebreak This case study extends invariant theory to computational crystallography, where a quantity struct \code{(struct quantity (value units) \#:transparent)} preserves dimensional invariants \code{(a . 1)(a . 1) = (a . 2)} and geometric relationships in TEM coordinate transformations. Here, \textit{Racket} evaluates trigonometric relationships using vector dimensional analysis. Similarly, Tsai and Melo (2014) compute their ``omni strain envelope'' by evaluating unit loading strain vectors from \(0\) to \(2\pi\): ``The inner envelope can be determined by finding the controlling ply that would fail first for unit loading strain vectors from 0 to \(2\pi\).'' This ``omni strain failure envelope'' constitutes an invariant that represents any laminate constructed with the material.
	
	To ensure the geometric integrity of molecular models, this paper develops a functional dimensional quantity system in \textit{Racket}, extends invariant theory to crystallographic trigonometry, and demonstrates how vector-based structural constraints can prevent dimensional errors in computational biology and chemistry.
	
	
	
	
	 
	
	

	\printbibliography
\end{document}